\section{Bridge II: Nonlinear Compactness}

The second bridge upgrades compactness to the level required for nonlinear closure. Once the cascade channel is controlled, the route still has to justify passage to the limit in the convective term without importing regularized structure. On the theorem surface this is achieved on one fixed classical approximation scheme.

\begin{proposition}[Classical nonlinear passage]
\label{prop:compactness}
There exists a classical approximation scheme \(u^{(n)}\to u\) strongly in \(L^2([0,T];L^2_{\mathrm{loc}}(\mathbb{R}^3))\) such that
\[
u^{(n)}\otimes u^{(n)} \to u\otimes u
\quad\text{in distributions on } (0,T)\times\mathbb{R}^3.
\]
Hence the convective term passes to the limit on the classical equation with no additional nonlinear defect term.
\end{proposition}

\begin{proof}
Let \(\varphi\in C_c^\infty((0,T)\times\mathbb{R}^3;\mathbb{R}^3)\) be divergence-free. By the fixed classical approximation scheme from the proposition,
\[
u^{(n)}\otimes u^{(n)} \to u\otimes u
\quad\text{in distributions on } (0,T)\times\mathbb{R}^3.
\]
Therefore
\[
\int_0^T\!\!\int_{\mathbb{R}^3}
\bigl(u^{(n)}\otimes u^{(n)}\bigr):\nabla\varphi\,dx\,dt
\to
\int_0^T\!\!\int_{\mathbb{R}^3}
\bigl(u\otimes u\bigr):\nabla\varphi\,dx\,dt.
\]
The strong local \(L^2\) convergence identifies the limiting field, while the tensor convergence excludes any extra defect contribution at the theorem level. This is exactly the nonlinear passage needed in the present route.
\end{proof}

The compactness step is therefore no longer a placeholder. It is the fixed nonlinear-closure package imported by the main theorem.
